\documentclass[11pt,a4paper]{article}
\usepackage[margin=26mm]{geometry}
\usepackage[T1]{fontenc}
\usepackage{lmodern,microtype,amsmath,amssymb,amsthm,mathtools}
\usepackage{xcolor,booktabs,enumitem}
\usepackage[colorlinks=true,linkcolor=blue!45!black,urlcolor=blue!50!black,citecolor=blue!45!black]{hyperref}
\hypersetup{pdftitle={A quartic entropy deficit for random rankings},pdfauthor={Dev D. Goyal}}
\newtheorem{theorem}{Theorem}
\newtheorem{lemma}[theorem]{Lemma}
\newtheorem{proposition}[theorem]{Proposition}
\newtheorem{remark}[theorem]{Remark}
\newcommand{\E}{\mathbb E}
\newcommand{\Prb}{\mathbb P}
\newcommand{\R}{\mathbb R}
\newcommand{\Sset}{\mathfrak S}
\newcommand{\eps}{\varepsilon}
\DeclareMathOperator{\Ent}{H}
\title{A quartic entropy deficit for random rankings}
\author{Dev D. Goyal}
\date{27 October 2026}
\setlength{\emergencystretch}{2em}
\begin{document}
\maketitle
\begin{abstract}
Let $P$ be any distribution on the rankings of an $n$-element set. If every
pairwise comparison has probability at distance at least $\eps>0$ from $1/2$,
then
\[
\Ent(P)\leq \left(1-\frac{\eps^4}{144}\right)\log(n!).
\]
The bound holds for every $n\geq0$, without a transitivity or independence
assumption. The proof combines the squared energy of the mean comparison
matrix, a normalized cut obtained from balanced splitting, and an exact
entropy bound for binary interleavings. Factorial induction removes any
large-size threshold. The accompanying Lean formalization uses finite probability distributions
and the actual comparison events of complete rankings.
\end{abstract}

\section{Statement and relation to prior work}

A ranking of a finite set $V$ is a list containing every element exactly once,
equivalently a bijection $V\to\{0,\ldots,n-1\}$. Write $u\prec v$ when $u$
occurs earlier. All logarithms are natural, and
\[
\Ent(P)=-\sum_{\rho\in\Sset_V} P(\rho)\log P(\rho),\qquad 0\log0:=0.
\]
Changing both entropy and logarithms to another fixed base greater than one leaves the
coefficient in the theorem unchanged.

\begin{theorem}\label{thm:main}
Let $V$ be a finite set of cardinality $n$, and let $P$ be a probability
distribution on its rankings. Suppose $\eps>0$ and
\begin{equation}\label{eq:bias}
\left|\Prb_P(u\prec v)-\frac12\right|\geq\eps
\quad\text{for every distinct }u,v\in V.
\end{equation}
Then
\begin{equation}\label{eq:main}
\Ent(P)\leq\left(1-\frac{\eps^4}{144}\right)\log(n!).
\end{equation}
\end{theorem}

For $n\leq1$ the ranking is unique and both sides vanish. For $n\geq2$, the
hypothesis itself forces $\eps\leq1/2$; an extra upper-bound assumption is
therefore unnecessary.

Leighton and Moitra introduced the entropy question in their work on poset
extensions and proved a quartic bound when the preferred comparisons are
consistent with a total order~\cite{LM}. Acan, Devlin, and Kahn proved the
general qualitative theorem~\cite{ADK}. Their general coefficient, at the
end of Section 3, has dependence $\exp[-\eps^{-O(1)}]$; their transitive
Theorem 2 gives $(1-\eps^2/8)n\log n$. Devlin's research statement explicitly
proposes a sufficient deficit of the form $C\eps^4$ in the general
log-factorial theorem~\cite{Devlin}.

Thus the qualitative local-to-global theorem is established prior work.
The contribution claimed here is the explicit general quartic bound and
the argument producing it. The constant and exponent are not claimed to be
optimal. 

\section{The mean comparison matrix and its energy}

Define
\[
B_{uv}=\Prb_P(u\prec v)-\Prb_P(v\prec u).
\]
Then $B_{uu}=0$, $B_{vu}=-B_{uv}$, and, for $u\ne v$,
$B_{uv}=2\Prb_P(u\prec v)-1$. Put $b=2\eps$. The hypothesis gives
$|B_{uv}|\geq b$ for distinct labels.

For an ordering $\tau=(v_1,\ldots,v_m)$ of some subset of $V$, define
\[
F_B(\tau)=\sum_{i<j}B_{v_i v_j}.
\]

\begin{lemma}\label{lem:energy}
For a random full ranking $\rho$ with law $P$,
\[
\E_P F_B(\rho)=\frac12\sum_{u,v\in V}B_{uv}^2.
\]
Consequently, if $n\geq2$ and all distinct entries have modulus at least $b$,
there is a ranking $\tau$ with $F_B(\tau)\geq b^2n^2/4$.
\end{lemma}
\begin{proof}
For each unordered pair $\{u,v\}$, its contribution to the score is $B_{uv}$
if $u$ precedes $v$, and $B_{vu}=-B_{uv}$ otherwise. Its expected contribution
is therefore
\[
B_{uv}\bigl(\Prb(u\prec v)-\Prb(v\prec u)\bigr)=B_{uv}^2.
\]
Summing over unordered pairs proves the identity. There are $n(n-1)$
ordered pairs of distinct elements, so the expectation is at least
$b^2n(n-1)/2\geq b^2n^2/4$. At least one value of a finite random variable
is no smaller than its expectation.
\end{proof}

\section{A normalized cut}

For disjoint subsets $L,R\subseteq V$, write
$B(L,R)=\sum_{u\in L,v\in R}B_{uv}$. When $n\geq2$, the following maximum
is over a finite nonempty family:
\[
A=\max_{\substack{L,R\ne\varnothing\\L\cap R=\varnothing}}
\frac{|B(L,R)|}{(|L|+|R|)^{3/2}}.
\]

\begin{lemma}\label{lem:ordering}
Every ordering $\tau$ of a subset of $V$, of length $s$, satisfies
$F_B(\tau)\leq6A s^{3/2}$.
\end{lemma}
\begin{proof}
Induct on $s$. The assertion is immediate for $s=0,1$. For $s\geq2$, split
the list into consecutive parts of lengths $a=\lfloor s/2\rfloor$ and
$r=\lceil s/2\rceil$. Both lengths are positive and at most $2s/3$.
The score is the sum of the two child scores and the cut sum between their
label sets. The latter is at most $A s^{3/2}$. Moreover,
\[
a^{3/2}+r^{3/2}
\leq(a+r)\sqrt{2s/3}
\leq\frac56s^{3/2},
\]
where the last inequality follows by squaring
$\sqrt{2/3}\leq5/6$. The induction hypotheses give a total at most
$6A(5/6)s^{3/2}+As^{3/2}=6As^{3/2}$.
\end{proof}

\begin{lemma}\label{lem:cut}
Under the hypotheses of Lemma~\ref{lem:energy}, there are disjoint nonempty
$L,R\subseteq V$ such that, with $q=|L|+|R|$ and $D=B(L,R)$,
\begin{equation}\label{eq:cut}
\frac{D^2}{q^3}\geq\frac{b^4n}{576}.
\end{equation}
\end{lemma}
\begin{proof}
Apply Lemma~\ref{lem:ordering} to the ordering from
Lemma~\ref{lem:energy}. It gives
$b^2n^2/4\leq6An^{3/2}$, hence $A\geq b^2\sqrt n/24$.
A pair attaining the maximum has $D^2/q^3=A^2$, which proves the claim.
The argument does not prescribe the sign of $D$.
\end{proof}

\section{Entropy of a binary interleaving}

Let $\Omega_{q,r}$ be the binary words of length $q$ containing exactly $r$
symbols $R$, and $q-r$ symbols $L$. Define
\begin{align*}
C(y)&=\#\{L\text{ before }R\}-\#\{R\text{ before }L\}\\
&=2\sum_{j=1}^{q}j\,\mathbf1_{\{y_j=R\}}-r(q+1).
\end{align*}
To verify the identity, the signed sum over all positions against an $R$
in position $j$ is $2j-q-1$. Removing the $R/R$ comparisons leaves the same
total, because those terms cancel in pairs. Equivalently, expand the
signed count over all pairs of positions and use antisymmetry.

\begin{lemma}\label{lem:two}
Suppose $\alpha,\beta\geq0$, $s\geq0$, $|u|,|v|\leq s$, and
$\alpha u+\beta v=0$. For every $t\in\R$,
\[
\alpha e^{tu}+\beta e^{tv}\leq(\alpha+\beta)e^{s^2t^2/2}.
\]
\end{lemma}
\begin{proof}
If $s=0$, both $u$ and $v$ vanish. Otherwise convexity of the exponential
on $[-s,s]$ gives, for $|x|\leq s$,
\[
e^{tx}\leq\frac{s+x}{2s}e^{ts}+\frac{s-x}{2s}e^{-ts}
=\cosh(ts)+\frac{x}{s}\sinh(ts).
\]
The weighted linear terms cancel. Finally
$\cosh z\leq e^{z^2/2}$: compare their even power series term by term using
$(2k)!=\prod_{j=1}^k(2j)(2j-1)\geq2^k k!$. All terms are nonnegative. This proves the bound for every
real $t$, including negative $t$.
\end{proof}

\begin{lemma}\label{lem:partition}
For every $q\geq0$, $0\leq r\leq q$, and $t\in\R$,
\begin{equation}\label{eq:partition}
Z_{q,r}(t):=\sum_{y\in\Omega_{q,r}}e^{tC(y)}
\leq\binom qr\exp(q^3t^2/2).
\end{equation}
\end{lemma}
\begin{proof}
Induct on $q$. The case $q=0$ has one empty word. For $q>0$, the extreme
slices $r=0,q$ each have one word, with $C=0$, and satisfy the estimate.
For $1\leq r<q$, splitting by the first symbol gives the exact recurrence
\[
Z_{q,r}(t)=e^{tr}Z_{q-1,r}(t)+e^{t(r-q)}Z_{q-1,r-1}(t).
\]
Indeed, a first $L$ precedes all $r$ later $R$ symbols, while a first $R$
precedes all $q-r$ later $L$ symbols and contributes $r-q$.
Set $\alpha=\binom{q-1}{r}$ and $\beta=\binom{q-1}{r-1}$. The binomial
identities
\[
\alpha+\beta=\binom qr,\qquad \alpha r+\beta(r-q)=0
\]
allow Lemma~\ref{lem:two} with $s=q$. The induction hypotheses imply
\[
Z_{q,r}(t)\leq\binom qr
\exp\!\left(\bigl((q-1)^3+q^2\bigr)t^2/2\right)
\leq\binom qr e^{q^3t^2/2}.
\]
The last step uses $(q-1)^3+q^2\leq q^3$. Impossible slices have zero
partition function and zero binomial coefficient; the formal statement
also includes them.
\end{proof}

\begin{lemma}\label{lem:wordentropy}
If $q>0$ and $Y$ has any distribution supported on $\Omega_{q,r}$, then,
with $D=\E C(Y)$,
\begin{equation}\label{eq:wordentropy}
\Ent(Y)\leq\log\binom qr-\frac{D^2}{2q^3}.
\end{equation}
For $q=0$, the entropy and the statistic are zero.
\end{lemma}
\begin{proof}
The finite Gibbs inequality has the following elementary proof.
For any finite support set $S$, probability masses $p_x$, and real numbers
$f_x$, let $Z=\sum_{x\in S}e^{f_x}$ and $Q_x=e^{f_x}/Z$. On the positive
support of $p$, apply $\log z\leq z-1$ to $z=Q_x/p_x$ and multiply by $p_x$.
The resulting sum is at most $\sum_xQ_x-\sum_xp_x=0$. Rearranging yields
\[
\Ent(p)\leq\log\sum_{x\in S}e^{f_x}-\sum_xp_x f_x.
\]
This argument permits zero probabilities and needs no positivity assumption
outside the support. Apply it with $f_y=tC(y)$ and then
Lemma~\ref{lem:partition}:
\[
\Ent(Y)\leq\log\binom qr+q^3t^2/2-tD.
\]
The admissible choice $t=D/q^3$, of either sign, gives the assertion.
\end{proof}

For a ranking of $L\cup R$, record only whether each successive element
lies in $L$ or $R$. Its word statistic is exactly the sum of signed
comparisons over $L\times R$. Therefore its expected statistic is
$D=B(L,R)$, regardless of dependencies among comparisons. Combining
Lemmas~\ref{lem:cut} and~\ref{lem:wordentropy} yields an interleaving entropy
loss of at least $b^4n/1152$.

\section{Factorial induction and proof of the theorem}

The signed-mean version is
\begin{equation}\label{eq:signed}
\Ent(P)\leq(1-c)\log(n!),\qquad c=b^4/2304,
\end{equation}
whenever every distinct entry of the mean comparison matrix has modulus
at least $b>0$. Proceed by strong induction on $n$. The cases $n=0,1$ have
one ranking and zero entropy.

For $n\geq2$, select $L,R$ by Lemma~\ref{lem:cut}. Let $Q=L\cup R$,
$W=V\setminus Q$, and write
\[
a=|L|,\quad r=|R|,\quad q=a+r,\quad w=|W|,\quad n=q+w.
\]
A full ranking is uniquely determined by its rankings on $L,R,W$, its
inner $L/R$ interleaving word $Y$, and its outer $Q/W$ interleaving word $Z$.
To reconstruct it, merge the $L$ and $R$ lists according to $Y$, then merge
the resulting list with the $W$ list according to $Z$. Thus the encoding is
injective. Entropy is unchanged by an injective encoding and is at most
the sum of the entropies of its coordinates. For completeness, an injective map retains each positive mass and only
adds zero masses outside its image, so its entropy is unchanged. For
subadditivity, compare a joint distribution to the product of its coordinate
marginals. That product is positive wherever the joint mass is positive.
The same $\log z\leq z-1$ argument as in Lemma~\ref{lem:wordentropy}
gives joint entropy at most the cross-entropy against that product.
Expanding the logarithm of the product yields the sum of marginal
entropies. Repeated application gives the required finite-coordinate bound.
The formalization proves these identities directly from finite sums.

The distributions on the three restricted lists are the unconditional
pushforwards of $P$. For two retained labels, restriction changes neither
their relative order nor its probability. Each restricted law therefore
satisfies the same bias condition. Since $a,r>0$, all three sizes $a,r,w$
are strictly smaller than $n$, even if $Q=V$ or $W$ is empty. The induction
hypotheses apply to every child.

By Lemma~\ref{lem:wordentropy} and the selected cut,
\[
\Ent(Y)\leq\log\binom qr-\frac{b^4n}{1152},
\qquad \Ent(Z)\leq\log\binom nq.
\]
For the second inequality it suffices to drop the nonnegative deficit in
the same lemma; alternatively, use the size of the outer word's support.
Consequently, on setting
\[
T=\log(a!)+\log(r!)+\log(w!),\qquad
M=\log\binom qr+\log\binom nq,
\]
this gives
\begin{equation}\label{eq:induct}
\Ent(P)\leq(1-c)T+M-2cn.
\end{equation}
Exact factorial cancellation gives $T+M=\log(n!)$. Moreover,
\[
M\leq(q+n)\log2\leq2n,
\]
since each binomial coefficient is at most the number of all binary words
of the same length, $q\leq n$, and $\log2\leq1$. As $c\geq0$, the right-hand
side of \eqref{eq:induct} is at most
\[
(1-c)(T+M)+cM-2cn\leq(1-c)\log(n!).
\]
This closes the induction. Finally $b=2\eps$ gives
$b^4/2304=\eps^4/144$, proving Theorem~\ref{thm:main}.\hfill$\square$

\section{Formal model and reproduction}

The Lean model represents a finite probability distribution by nonnegative
real masses whose sum is one. Entropy is the finite sum of
\texttt{Real.negMulLog}. A ranking is a list with no repeated labels that
contains every label. Restriction filters this list. Pushforward probabilities
sum the masses over each inverse image. Comparison probabilities are
expectations of event indicators.

The main theorem is
\begin{center}\small\texttt{PermEntropy.FinDist.entropy\_of\_pairwise\_bias\_all}.\end{center}
The theorem \texttt{entropy\_of\_pairwise\_bias\_equiv\_all} gives the same
bound for bijections from the label set to its positions. The event in its
hypothesis is the literal comparison of two positions.
Both statements include every finite size. The central entropy and cut
bounds are proved lemmas, not assumptions of either theorem.

The package pins Lean 4.19.0 and the Mathlib revision in
\texttt{lake-manifest.json}. Its README gives the build and axiom-audit
commands. The audit permits only \texttt{propext},
\texttt{Classical.choice}, and \texttt{Quot.sound}.

\section{Scope and limitations}

The hypothesis concerns every unordered pair, with a common positive lower
bound on absolute bias. The preferred directions may contain cycles. The
argument places no independence condition on the distribution, on
restricted rankings, or on interleaving words. Its recursion uses
unconditional restrictions throughout.

The cut maximum is used existentially. This proof does not provide a
polynomial-time algorithm for finding the cut from a succinctly represented
distribution, nor an algorithm for realizing a prescribed comparison
matrix. It also does not determine the best exponent or constant, classify
extremizers, or prove a matching lower construction. The uniform
distribution is excluded by positive pairwise bias when $n\geq2$.

The balanced ordering decomposition has antecedents in tournament-ranking
arguments acknowledged in the cited literature. Gibbs' inequality, entropy
subadditivity, and the exponential estimate are standard tools. The
quartic bound combines these tools with the mean comparison energy,
normalized cuts, and exact factorial induction.

\section*{Autonomous AI research}
This work is fully autonomous agentic research. The initial research and proofs were generated autonomously by GPT 6 with Pro Reasoning in ChatGPT. GPT 6 Astra in Codex autonomously produced the formalizations and corrections. The authors checked the proofs to the best of their mathematical abilities. This check is not expert validation or peer review. Dev D. Goyal is responsible for publication.

\begin{thebibliography}{9}
\bibitem{LM}
T. Leighton and A. Moitra.
\emph{On Entropy and Extensions of Posets}. Manuscript, 12 September 2011.
\url{https://people.csail.mit.edu/moitra/docs/poset.pdf}.

\bibitem{ADK}
H. Acan, P. Devlin, and J. Kahn.
Proof of an entropy conjecture of Leighton and Moitra.
\emph{Journal of Combinatorial Theory, Series A} \textbf{161} (2019), 299--308.
\href{https://doi.org/10.1016/j.jcta.2018.08.004}{doi:10.1016/j.jcta.2018.08.004}.
\href{https://arxiv.org/abs/1701.04321}{arXiv:1701.04321v2}, 10 March 2017.

\bibitem{Devlin}
P. Devlin. \emph{Research Statement: Unabridged}. Undated primary document,
pp. 1 and 5. Accessed 27 October 2026.
\url{https://sites.math.rutgers.edu/~prd41/employers/Research_statement_unabridged.pdf}.

\bibitem{Lean}
The Lean and Mathlib communities. Lean 4 and Mathlib source repositories.
\url{https://github.com/leanprover/lean4};
\url{https://github.com/leanprover-community/mathlib4}.
Exact revisions appear in \texttt{lean-toolchain} and \texttt{lake-manifest.json}.
\end{thebibliography}
\end{document}
